\documentclass{article}
\usepackage{graphicx} % Required for inserting images
\usepackage{xcolor}
\usepackage{booktabs} 
\usepackage{multirow}
\usepackage{makecell}
\usepackage{float}
\usepackage{hyperref}
\hypersetup{
    colorlinks,
    linkcolor=blue
}
\title{Shifumi de Nash}
\author{Pierre Nunn et Léo Laffeach}

\begin{document}

\maketitle
\tableofcontents
\newpage
\section{Objectifs pédagogiques}

Le but de cette activité est de faire comprendre aux élèves les notions d'équilibre de Nash et d'Optimum de Pareto en utilisant le jeu pierre-feuille-ciseaux ou shifumi. (Ne pas forcément utiliser les termes avant le "c'est de l'informatique parce que \dots")

\section{Définitions}

En théorie des jeux, en considérant un jeu à plusieurs joueurs, on cherche l'ensemble des stratégies de chaque joueur pour maximiser leur gain. 

Une stratégie est un ensemble de choix durant la partie. Elle peut être dépendante ou indépendante de la stratégie des autres joueurs.
\\
\\
Optimum de Pareto : 

Un optimum de Pareto est une stratégie pour laquelle il n'y a pas de moyen d'augmenter les gains d'un joueur sans diminuer les gains d'un autre. 

C'est-à-dire, si on note les stratégies $s_1,..., s_n$, les gains du joueur i $g_i(s_1, ..., s_n)$ et $(s*_i)_{1\leq i \leq n}$ l'optimum de Pareto, on a que s'il existe un joueur i et des stratégies $(s'_i)_{1\leq i \leq n}$ telles que $g_i(s'_1, ..., s'_n) > g_i(s*_1, ..., s*_n)$ alors, il existe un joueur j tel que $g_j(s'_1, ..., s'_n) < g_j(s*_1, ..., s*_n)$.

Il peut y avoir plusieurs optima de Pareto.
\\
\\
Optimum social :

Un optimum social est un optimum de Pareto particulier, pour lequel la somme des gains de tous les joueurs est optimale.

Ce n'est pas forcément un équilibre de Nash. (Dilemme du prisonnier)
\\
\\
Équilibre de Nash : 

Un équilibre de Nash est une stratégie qui maximise le gain de tous les joueurs.

C'est-à-dire que pour tout joueur i, $g_i(s*_1, ..., s_i, ..., s*_n) \leq g_i(s*_1, ..., s*_i, ..., s*_n)$

Il n'en existe pas toujours un.

\section{Matériels}
Des cartes avec un symbole de pierre, de feuille et de ciseaux. \\
Assez de cartes pour que tout le monde ait 6 cartes (2 de chaques type).
\newpage
\section{Activité}

Le principe de l'activité est que les élèves jouent des parties en un nombre fixé $n$ de manches de pierre-feuille-ciseaux. Avant qu'ils ne commencent une partie, un coup (pierre, feuille ou ciseaux) est annoncé à l'avance (ça peut être le même coup pour toutes les parties ou ça peut se décider avec un dé à chaque fois). Les élèves gagneront des points selon le fait qu'ils aient joué ce qui est annoncé ou selon leur résultat face à ce que l'autre élève a joué.

On dit aux élèves que le but du jeu est d'avoir le plus de points possibles à la fin des $n$ manches.Certains voudront probablement essayer de battre leur adversaire en ayant plus de points que lui, mais ce n'est pas une stratégie optimale pour gagner des points. Lors d'une remise en commun, on comparera les stratégies de chaque paire de joueurs et on constatera (normalement) que ce sont les joueurs qui ont collaboré plutôt que de chercher à battre l'adversaire qui ont obtenu le plus de points et que dans ce jeu, les joueurs ont un intérêt à s'entraider. (On peut amener à ce moment-là les notions d'équilibre de Nash et d'optimum de Pareto).
\\
\\
On peut alors compter les points ainsi:
\begin{itemize}
    \item Victoire : +2 points
    \item Égalité : +1 points
    \item Défaite : +0 points
    \item Annoncé : +1 points
\end{itemize}
On peut aussi faire varier les gains pour en faire un jeu à somme nulle mais ça modifie les propriétés du jeu. 

\section{Déroulement de l'activité}
On pourra utiliser des cartes pour faciliter les parties. Chaque joueur a 6 cartes, 2 de chaque type.

\subsection{Découverte}
Faire jouer les élèves au pierre-papier-ciseaux classique.\\
Par exemple, demander aux élèves de se mettre par deux et de faire 5-10 parties, en comptant les points.
\subsection{$1^{ère}$ mise en commun}
Demander combien de points ont eu certains élèves et s'ils avaient une stratégie. 
Il peut y avoir des stratégies entre parties, comme jouer la même chose qu'avant ou jouer le faible de ce que l'autre a joué avant. Mais bien faire comprendre qu'on parle de stratégie pour la partie en cours sans considérer les autres parties.\\
Il n'y a pas vraiment de stratégie, on ne sait pas ce que l'autre va jouer donc on joue un peu à l'aveugle. En moyenne, tout le monde gagne le même nombre de points.
\\
\\
Introduire la variante de l'annoncé. Au début de chaque partie, on tire une carte qui détermine quel coup est "meilleur". Si les joueurs jouent ce coup, ils gagnent un point. 

\begin{table}[H]
\renewcommand{\arraystretch}{2}%
\centering
\aboverulesep = 0pt
\belowrulesep = 0pt
\begin{tabular}{cc|c|c|c}
          & & \multicolumn{3}{c}{J2} \\ 
    & J1/J2 & faible & annoncé & fort \\ \midrule
    & faible & 1/1 & \textcolor{red}{0/3} & 2/0  \\ \cmidrule{2-5}
    J1 &annoncé & \textcolor{red}{3/0} & \textcolor{purple}{2/2} & 1/2 \\ \cmidrule{2-5}
     & fort & 0/2 & 2/1 & 1/1 
\end{tabular}
\end{table}

\begin{itemize}
    \item Bleu : Équilibre de Nash
    \item Rouge : Optimum de Pareto
    \item Violet : Les deux
\end{itemize}

\subsection{Jouer avec l'annonce}
On laisse les élèves jouer selon cette variante 5-10 parties, en comptant les points.

\subsection{$2^{ème}$ mise en commun}
On demande les mêmes choses, combien de points et s'ils ont une stratégie.\\
Dessiner le tableau des gains sur un exemple et montrer les optima de Pareto. On pourra supposer que le jeu n'est pas exactement simultané pour en parler. Présenter l'équilibre de Nash, la stratégie qui maximise le gain des deux joueurs, si on additionne leurs scores.
\\
\\
Introduire une variante: 
\subsubsection{estimation de partie}
    Les deux joueurs tirent une carte avant la partie. Si un joueur joue ce que sa carte indique, il gagne un point. 

    Cet extension laisse les joueurs annoncer simultanément un coup. Elle n'est pas forcément intéressante en pratique : elle est difficile et ne permet pas de voir clairement la notion d'équilibre de Nash (sauf l'équilibre de Nash probabiliste mais quitte à en parler, autant le faire sur le jeu de Pierre-feuille-ciseaux normal sans annonce).


    \begin{table}[H]
    \renewcommand{\arraystretch}{2}%
    \centering
    \aboverulesep = 0pt
    \belowrulesep = 0pt
        \begin{tabular}{|c|c||c|c|c||c|c|c||c|c|c|}
            \toprule
            J1/J2  & J1 &  \multicolumn{3}{|c||}{Pierre} & \multicolumn{3}{c||}{Feuille} & \multicolumn{3}{c|}{Ciseaux}   \\ 
            \cmidrule{1-11}
              J2   &   & P &   F    & C & P &    F    & C & P &    F    & C \\ \midrule \midrule
                  & P &\textcolor{red}{2/2}&  2/1   & 0/3 & 1/2 &  \textcolor{red}{3/1}  & 0/3 & 1/2 &   2/1   &\textcolor{red}{1/3}\\ 
            \cmidrule{2-11}
             Pierre & F & 1/2 &  1/1   & 2/0 & 0/2 &   2/1   & 2/0 & 0/2 &   1/1   & 3/0 \\ 
            \cmidrule{2-11}
                   & C & 3/0 &  0/2   & 1/1 & 2/0 &   1/2   & 1/1 & 2/0 &   0/2   & 2/1 \\ \midrule \midrule
                   & P & 2/1 &  2/0   & 0/2 & 1/1 &   3/0   & 0/2 & 1/1 &   2/0   & 1/2 \\ 
            \cmidrule{2-11}
            Feuille& F &\textcolor{red}{1/3}&  1/2   & 2/1 & 0/3 &  \textcolor{red}{2/2}  & 2/1 & 0/3 &   1/2   &\textcolor{red}{3/1}\\ 
            \cmidrule{2-11}
                   & C & 3/0 &  0/2   & 1/1 & 2/0 &   1/2   & 1/1 & 2/0 &   0/2   & 2/1 \\ \midrule \midrule
                   & P & 2/1 &  2/0   & 0/2 & 1/1 &   3/0   & 0/2 & 1/1 &   2/0   & 1/2 \\ 
            \cmidrule(l){2-11}
            Ciseaux& F & 1/2 &  1/1   & 2/0 & 0/2 &   2/1   & 2/0 & 0/2 &   1/1   & 3/0\\ 
            \cmidrule{2-11}
                   & C &\textcolor{red}{3/1}& 0/3  & 1/2 & 2/1 &  \textcolor{red}{1/3}  & 1/2 & 2/1 &  0/3  &\textcolor{red}{2/2}\\ \bottomrule
        \end{tabular}
        \caption{Score selon l'annoncé par rapport au joué pour chaque joueur}
    \end{table}

\subsubsection{Dilemme du prisonnier}

    Jouer l'annonce donne 2 points à l'autre, ne pas le faire fait gagner 1 points, en plus des points de la partie. 

    \begin{itemize}
        \item Jouer l'annoncé : Donner 2 points à l'autre
        \item Ne pas le jouer : Gagner un point
    \end{itemize}
    
    \begin{table}[H]
    \renewcommand{\arraystretch}{2}%
    \centering
    \aboverulesep = 0pt
    \belowrulesep = 0pt
    \begin{tabular}{c|c|c}
        J1/J2 & annoncé & autre  \\ \midrule
        annoncé & \textcolor{red}{2/2} & \textcolor{red}{0/3}  \\ \midrule
        autre & \textcolor{red}{3/0} & \textcolor{blue}{1/1} \\ 
    \end{tabular}
    \end{table}

\subsubsection{Chicken}

    Jouer autre chose que l'annonce fait gagner 1 points, mais si un seul joueur joue l'annonce, il gagne 4 points, en plus des points de la partie.
    
    Pénaliser le cas d'égalité dans l'annonce.
    Si seulement un des joueurs change, qu'il soit gagnant ou perdant, les joueurs gagnent des points. Mais si les deux changent, les deux perdent des points. L'intérêt reste d'avoir le plus de points sans se concerter.
    
    \begin{table}[H]
    \renewcommand{\arraystretch}{2}%
    \centering
    \aboverulesep = 0pt
    \belowrulesep = 0pt
    \begin{tabular}{c|c|c}
        J1/J2 & annoncé & autre  \\ \midrule
        annoncé & {0/0} & \textcolor{red}{4/1}  \\ \midrule
        autre & \textcolor{red}{1/4} & \textcolor{blue}{1/1} \\ 
    \end{tabular}
    \end{table}
\subsection{Jouer avec la variante}
On laisse les élèves jouer selon cette variante 4-5 parties, en comptant les points. 
\\
\\
Selon le temps, on peut évoquer d'autres variantes.

\subsection{Dernière mise en commun}
Nombre de points et existence d'une stratégie ? \\
Demander si c'est différent d'une partie de base ? \\
Expliquer que tous les jeux n'ont pas forcément de stratégie menant à un équilibre de Nash et que bien que certains jeux semblent différents, ils peuvent avoir les mêmes stratégies et résultats.\\
Expliquer en quoi c'est de l'informatique.

\section{Exemples informatiques}
La théorie des jeux est utilisée en informatique pour déterminer des stratégies optimales lors de l'accès à des ressources limitées.
\begin{itemize}
    \item TCP
    \item accessibilité à internet (wifi, 3G, ADSL...)
    \item équité entre les participants
\end{itemize}

\end{document}